\documentclass[12pt,letterpaper,oneside]{extreport}

% Reporte de funcionamiento de la encuesta FES Aragon
\usepackage[utf8]{inputenc}
\usepackage[T1]{fontenc}
\usepackage[spanish,es-nodecimaldot]{babel}
\usepackage{geometry}
\geometry{letterpaper,left=3cm,right=3cm,top=2.5cm,bottom=2.5cm}
\usepackage{setspace}
\onehalfspacing
\usepackage{graphicx}
\usepackage{xcolor}
\usepackage{hyperref}
\usepackage{booktabs}
\usepackage{longtable}
\usepackage{float}
\usepackage{enumitem}
\usepackage{listings}
\usepackage{tikz}
\usetikzlibrary{arrows.meta,positioning,shapes.geometric}
\usepackage{fancyhdr}
\usepackage{titlesec}
\usepackage{parskip}

\definecolor{azul}{HTML}{173F5F}
\definecolor{azulClaro}{HTML}{20639B}
\definecolor{amarillo}{HTML}{F6C85F}
\definecolor{grisCodigo}{HTML}{F3F4F6}
\definecolor{grisTexto}{HTML}{475569}

\hypersetup{colorlinks=true,linkcolor=azul,urlcolor=azulClaro}
\pagestyle{fancy}
\fancyhf{}
\lhead{\small Encuesta de Experiencia Universitaria}
\rhead{\small FES Aragon}
\cfoot{\thepage}
\renewcommand{\headrulewidth}{0.4pt}

\titleformat{\chapter}{\Huge\bfseries\color{azul}}{\thechapter.}{1em}{}
\titleformat{\section}{\Large\bfseries\color{azul}}{\thesection}{1em}{}
\titleformat{\subsection}{\large\bfseries\color{azulClaro}}{\thesubsection}{1em}{}
\titlespacing*{\chapter}{0pt}{-25pt}{20pt}

\lstset{
  basicstyle=\ttfamily\small,
  backgroundcolor=\color{grisCodigo},
  frame=single,
  breaklines=true,
  showstringspaces=false,
  columns=fullflexible
}

\begin{document}

\begin{titlepage}
\centering
\vspace*{1.2cm}
{\Huge\bfseries Universidad Nacional Autónoma de México\par}
\vspace{0.5cm}
{\LARGE Facultad de Estudios Superiores Aragón\par}
\vspace{2cm}
\rule{\linewidth}{1.5pt}
\vspace{0.7cm}
{\Huge\bfseries Reporte de funcionamiento\par}
\vspace{0.4cm}
{\LARGE Encuesta de Experiencia Universitaria\par}
\vspace{0.7cm}
\rule{\linewidth}{1.5pt}
\vfill
\begin{flushleft}
\large
\textbf{Asignatura:} Desarrollo Web\\[0.35cm]
\textbf{Proyecto:} Aplicación web de encuesta y análisis de respuestas\\[0.35cm]
\textbf{Tecnologías:} PHP, MySQL, HTML, CSS, JavaScript, Chart.js y Docker\\[0.35cm]
\textbf{Autor:} [Escribe tu nombre]\\[0.35cm]
\textbf{Fecha:} \today
\end{flushleft}
\end{titlepage}

\tableofcontents
\newpage

\chapter{Resumen ejecutivo}

El presente documento describe el funcionamiento de una aplicación web diseñada para capturar y analizar una encuesta de experiencia universitaria. El sistema permite que una persona entre directamente a la encuesta, registre sus datos básicos y responda seis preguntas cerradas y dos preguntas abiertas.

La información se almacena en una base de datos MySQL. Después de guardar la encuesta, la aplicación muestra un resumen individual y gráficas interactivas. En la página principal también se presentan las distribuciones acumuladas de las respuestas almacenadas, incluyendo los registros iniciales y las nuevas respuestas.

La aplicación utiliza PHP para la lógica del servidor, MySQL para persistencia, HTML y CSS para la interfaz, JavaScript con Chart.js para las gráficas, y Docker Compose para levantar los servicios web y de base de datos.

\chapter{Objetivos}

\section{Objetivo general}

Desarrollar una aplicación web capaz de capturar respuestas de una encuesta y mostrar información estadística actualizada a partir de los registros almacenados en MySQL.

\section{Objetivos específicos}

\begin{enumerate}[label=\arabic*.]
    \item Registrar sexo y edad de cada participante.
    \item Validar que la edad mínima para responder sea de 20 años.
    \item Capturar seis preguntas cerradas con cuatro opciones cada una.
    \item Capturar dos respuestas abiertas.
    \item Guardar cada interacción en la tabla \texttt{encuesta\_respuestas}.
    \item Calcular frecuencias y porcentajes con base en todos los registros.
    \item Presentar gráficas circulares y demográficas de sexo y edad.
\end{enumerate}

\chapter{Descripción general del sistema}

La aplicación sigue un flujo sencillo para reducir pasos innecesarios para el participante:

\begin{enumerate}[label=\arabic*.]
    \item La persona entra a la URL de la aplicación.
    \item El sistema crea o reutiliza una sesión anónima.
    \item La aplicación muestra directamente la encuesta.
    \item El servidor valida todos los datos recibidos.
    \item La respuesta se inserta en MySQL.
    \item El usuario es enviado a la página de resumen.
    \item La aplicación consulta todos los registros y genera las gráficas.
\end{enumerate}

\begin{figure}[H]
\centering
\begin{tikzpicture}[
  node distance=1.2cm,
  bloque/.style={draw, rounded corners, fill=blue!8, align=center, minimum width=5.2cm, minimum height=0.85cm},
  flecha/.style={-{Latex}, thick, color=azul}
]
\node[bloque] (inicio) {Entrada automática};
\node[bloque, below=of inicio] (encuesta) {Formulario de encuesta};
\node[bloque, below=of encuesta] (bd) {MySQL: encuesta\_respuestas};
\node[bloque, below=of bd] (graficas) {Resumen y gráficas interactivas};
\draw[flecha] (inicio) -- (encuesta);
\draw[flecha] (encuesta) -- (bd);
\draw[flecha] (bd) -- (graficas);
\end{tikzpicture}
\caption{Flujo general de la aplicación}
\end{figure}

\chapter{Cuestionario implementado}

La encuesta contiene dos datos demográficos, seis preguntas cerradas y dos preguntas abiertas. La edad se valida en el servidor y en el navegador con un rango de 20 a 99 años.

\begin{longtable}{p{0.08\textwidth}p{0.42\textwidth}p{0.38\textwidth}}
\caption{Preguntas cerradas y opciones}\label{tab:preguntas}\\
\toprule
\textbf{No.} & \textbf{Pregunta} & \textbf{Opciones}\
\midrule
\endfirsthead
\toprule
\textbf{No.} & \textbf{Pregunta} & \textbf{Opciones}\
\midrule
\endhead
1 & Satisfacción con la limpieza del campus & Muy satisfecho; Satisfecho; Poco satisfecho; Nada satisfecho\\
2 & Espacio utilizado con mayor frecuencia & Aulas; Biblioteca; Áreas deportivas; Laboratorios\\
3 & Medio para enterarse de actividades universitarias & Redes sociales; Carteles; Docentes; Compañeros\\
4 & Aspecto que se mejoraría primero en la FES & Seguridad; Conectividad; Áreas verdes; Servicios escolares\\
5 & Frecuencia de participación en actividades culturales & Cada semana; Cada mes; Ocasionalmente; Nunca\\
6 & Horario preferido para actividades extracurriculares & Antes de clases; Entre clases; Después de clases; Fin de semana\\
\bottomrule
\end{longtable}

\section{Preguntas abiertas}

\begin{enumerate}[label=\arabic*.]
    \item ¿Qué cambio concreto haría más agradable tu experiencia en el campus?
    \item ¿Qué actividad o servicio nuevo te gustaría encontrar en la FES Aragón?
\end{enumerate}

\chapter{Base de datos}

La tabla principal de respuestas es \texttt{encuesta\_respuestas}. Cada fila representa una interacción completa con la encuesta.

\begin{longtable}{p{0.28\textwidth}p{0.58\textwidth}}
\caption{Campos de la tabla de respuestas}\label{tab:bd}\\
\toprule
\textbf{Campo} & \textbf{Descripción}\
\midrule
\endfirsthead
\toprule
\textbf{Campo} & \textbf{Descripción}\
\midrule
\endhead
\texttt{id} & Identificador autoincremental de la respuesta.\\
\texttt{user\_id} & Identificador del usuario asociado a la sesión.\\
\texttt{sexo} & Categoría demográfica: Mujer, Hombre, No binario o Prefiero no decirlo.\\
\texttt{edad} & Edad de la persona, validada entre 20 y 99 años.\\
\texttt{pregunta\_1} a \texttt{pregunta\_6} & Respuestas cerradas del cuestionario.\\
\texttt{abierta\_1}, \texttt{abierta\_2} & Respuestas de texto libre.\\
\texttt{created\_at} & Fecha y hora de creación del registro.\\
\bottomrule
\end{longtable}

Para la entrega inicial se preparan 98 registros. La aplicación puede recibir posteriormente dos registros adicionales para alcanzar 100 respuestas.

\chapter{Funcionamiento de las gráficas}

\section{Gráficas de preguntas cerradas}

La página de encuesta consulta la tabla de respuestas y cuenta cuántas veces aparece cada opción válida. Los resultados se envían al navegador en formato JSON y Chart.js genera una gráfica circular para cada pregunta.

El porcentaje de una opción se calcula mediante:

\[
\text{Porcentaje} = \frac{\text{frecuencia de la opción}}{\text{total de respuestas}} \times 100
\]

Cada pregunta conserva sus cuatro opciones y la suma de sus frecuencias representa el total de registros válidos.

\section{Gráfica de edad}

La edad se agrupa en cuatro rangos:

\begin{itemize}
    \item 20 a 24 años.
    \item 25 a 29 años.
    \item 30 a 39 años.
    \item 40 años o más.
\end{itemize}

\section{Gráfica de sexo}

La distribución por sexo utiliza las categorías almacenadas en la base de datos: Mujer, Hombre, No binario y Prefiero no decirlo.

Las gráficas incluyen información emergente o \textit{tooltip} con la frecuencia y el porcentaje correspondiente.

\section{Historial de respuestas abiertas}

El sistema incorpora un apartado administrativo denominado ``Historial de respuestas abiertas''. Esta vista consulta los campos \texttt{abierta\_1}, \texttt{abierta\_2} y \texttt{created\_at} de la tabla \texttt{encuesta\_respuestas}. Las entradas se presentan en orden cronológico descendente para que la administración pueda revisar primero las respuestas más recientes.

Cada entrada incluye el identificador de la respuesta, la fecha de registro, sexo, edad y el contenido de ambas preguntas abiertas. El acceso al historial está restringido a usuarios con rol administrativo.


\chapter{Arquitectura y despliegue}

\section{Tecnologías}

\begin{itemize}
    \item \textbf{PHP 8.2 y Apache:} procesamiento del servidor y páginas dinámicas.
    \item \textbf{MySQL 8:} almacenamiento de respuestas.
    \item \textbf{HTML5 y CSS3:} estructura y diseño responsive.
    \item \textbf{JavaScript y Chart.js:} gráficas interactivas.
    \item \textbf{Docker Compose:} orquestación de los servicios web y MySQL.
\end{itemize}

\section{Servicios de Docker}

\begin{table}[H]
\centering
\caption{Servicios definidos en Docker Compose}
\begin{tabular}{lll}
\toprule
\textbf{Servicio} & \textbf{Contenedor} & \textbf{Puerto}\
\midrule
Web & \texttt{fes\_web} & 8080:80\\
Base de datos & \texttt{fes\_mysql} & 3307:3306\\
\bottomrule
\end{tabular}
\end{table}

MySQL incluye un \textit{healthcheck}. El servicio web depende de que la base de datos se encuentre saludable antes de iniciar, lo que evita errores de conexión durante el arranque.

\section{Comandos de ejecución}

\begin{lstlisting}[language=bash]
docker compose up -d --build
docker compose ps
docker compose logs --tail=100
\end{lstlisting}

Para cargar los registros iniciales en una base existente:

\begin{lstlisting}[language=bash]
Get-Content database/restore_hobbies.sql |
  docker compose exec -T mysql mysql -uroot -proot -Dfesaragon
\end{lstlisting}

\chapter{Validaciones y pruebas}

\begin{longtable}{p{0.35\textwidth}p{0.5\textwidth}}
\caption{Pruebas funcionales realizadas}\label{tab:pruebas}\\
\toprule
\textbf{Prueba} & \textbf{Resultado esperado}\
\midrule
\endfirsthead
\toprule
\textbf{Prueba} & \textbf{Resultado esperado}\
\midrule
\endhead
Acceso a la raíz & Redirección automática a la encuesta.\\
Edad menor a 20 & El servidor rechaza la respuesta.\\
Formulario incompleto & Se muestran mensajes de validación.\\
Envío válido & Se inserta una fila en MySQL.\\
Nueva respuesta & Aumenta el total mostrado por las gráficas.\\
Gráfica por edad & Agrupa los registros en cuatro rangos.\\
Gráfica por sexo & Cuenta las categorías guardadas.\\
Error de MySQL al iniciar & Docker espera el estado saludable del servicio.\\
\bottomrule
\end{longtable}

\chapter{Conclusiones}

La aplicación cumple con el flujo requerido para capturar, guardar y analizar respuestas de una encuesta. La información se conserva en MySQL y las gráficas se actualizan con base en todos los registros disponibles, incluyendo los datos iniciales y las nuevas interacciones.

La separación entre formulario, persistencia y visualización facilita el mantenimiento del proyecto. Además, Docker Compose permite levantar el entorno de desarrollo de manera reproducible y el diseño responsive permite utilizar la encuesta desde distintos tamaños de pantalla.

Como trabajo futuro se pueden incorporar exportación de resultados a PDF, filtros por fecha y una vista administrativa con más comparaciones entre variables demográficas.

\end{document}
